\subsection{A motivating example}

Suppose we want to predict how people will share resources in a novel public goods game (as they do in \cite{alsobay2025integrative}).
In this new game, for which we have no previous data, three participants will be endowed with \$5 and can choose to contribute any portion of their endowment, which will be multiplied by 3 and then divided equally among all participants.
While we do not have any prior public goods game data, we do have human data from a related \$20 dictator game.
It is related in the sense that one might reasonably expect some generalizable features of human choice to affect allocations in both games.
The observed human offers from the dictator game are \{6,6,7,7,8,8,9,9\}.

Before LLMs, one might have tried to train a standard machine learning model (e.g., decision trees or regression) solely on these dictator game offers.
However, such models would struggle to transfer to the structurally distinct public goods game, as they cannot adapt across different game formats without retraining.
LLMs offer a more flexible alternative. 
Rather than training a new model from scratch, we can prompt 
an existing LLM to simulate responses by instructing it to behave as a human participant.
For instance, we might use a baseline system prompt: \emph{``You are a human.''}
We can then append game-specific prompts such as \emph{``You are playing a dictator game with \$20...''} or \emph{``You are playing a public goods game with \$5...''} without any additional training.

Suppose that, when prompted with the dictator game instructions, the LLM produces a response distribution \{3,3,3,3,3,4,4,4\}.
Although well within the allowable offers, this distribution clearly diverges from the observed human data  (\{6,6,7,7,8,8,9,9\}) and leaves us with little hope that it could effectively predict responses to the novel public goods game.
Thus, even though the model may capture some general aspects of human behavior as a baseline, i.e., a tendency to offer nontrivial amounts \citep{Henrich2001InSearch}, it does so imperfectly.

One might think this problem could be addressed by first randomly splitting the human sample of dictator game offers into training and testing sets.
And then identifying the best-performing \persona{s} on the training set, and validating their performance on the test set \citep{ML2017Sendhil,ludwig2025llmapplied}.
Indeed, an LLM instructed \emph{``You randomly choose numbers between 6 and 9''} could reasonably predict both training and testing sets for any split of the human data better than the baseline LLM with no persona.
However, such an approach will almost certainly fail to generalize beyond the specific data-generating process that produced the dictator game offers.
It has, in effect, still overfit.
Rather than overfitting to the training data, it has overfit to the whole data-generating process.

Now, consider a more theoretically grounded \persona{}: \emph{``You are self-interested but fair.''}
Suppose an LLM endowed with this \persona{} also generates offers in the 6-9 range for the \$20 dictator game.
Crucially, this \persona{} aligns with known causal factors that drive human sharing behavior more broadly, like behavioral parameters in empirically tested theoretical models (e.g., \cite{charness2002understanding}).
It is a flexible decision-making program that can be effectively applied to various types of allocation games.
Yet, without explicit knowledge of the causal drivers governing behavior in the new public goods game, nothing in the data produced by either \persona{} alone rules out the atheoretical random-number prompt.
This illustrates the core challenge.
We seek to construct and select \persona{s} that do not overfit to a single data-generating process, but instead capture the stable behavioral drivers relevant across settings.
Then, we might gain confidence that they will better predict human responses in novel target settings governed by the same drivers.
